\documentclass{article}

% Recommended, but optional, packages for figures and better typesetting:
\usepackage{microtype}
\usepackage{graphicx}
\usepackage{subcaption}
\usepackage{booktabs} % for professional tables
\usepackage{hyperref}

% Attempt to make hyperref and algorithmic work together better:
\newcommand{\theHalgorithm}{\arabic{algorithm}}

% Use the following line for the initial blind version submitted for review:
\usepackage{icml2026}

\usepackage{amsmath}
\usepackage{amssymb}
\usepackage{mathtools}
\usepackage{amsthm}
\usepackage{algorithm}
\usepackage{algorithmic}

% if you use cleveref..
\usepackage[capitalize,noabbrev]{cleveref}
\usepackage{enumitem}
\setlist{nosep}

%%%%%%%%%%%%%%%%%%%%%%%%%%%%%%%%
% THEOREMS
%%%%%%%%%%%%%%%%%%%%%%%%%%%%%%%%
\theoremstyle{plain}
\newtheorem{theorem}{Theorem}[section]
\newtheorem{proposition}[theorem]{Proposition}
\newtheorem{lemma}[theorem]{Lemma}
\newtheorem{corollary}[theorem]{Corollary}
\theoremstyle{definition}
\newtheorem{definition}[theorem]{Definition}
\newtheorem{assumption}[theorem]{Assumption}
\theoremstyle{remark}
\newtheorem{remark}[theorem]{Remark}

\icmltitlerunning{Fisher-Weighted Procrustes Alignment for Compatible Model Merging}

\begin{document}

\twocolumn[
  \icmltitle{Fisher-Weighted Procrustes Alignment for Compatible Model Merging}

  \icmlsetsymbol{equal}{*}

  \begin{icmlauthorlist}
    \icmlauthor{Firstname1 Lastname1}{equal,yyy}
    \icmlauthor{Firstname2 Lastname2}{equal,comp}
    \icmlauthor{Firstname3 Lastname3}{sch}
  \end{icmlauthorlist}

  \icmlaffiliation{yyy}{Department of Computer Science, University of Research, Location, Country}
  \icmlaffiliation{comp}{AI Research Lab, Tech Company, Location, Country}
  \icmlaffiliation{sch}{School of Engineering, Institute of Technology, Location, Country}

  \icmlcorrespondingauthor{Firstname1 Lastname1}{first1.last1@research.edu}

  \icmlkeywords{Model Merging, Procrustes Alignment, Fisher Information, Out-of-Distribution Robustness, Lie Algebra}

  \vskip 0.3in
]

\begin{abstract}
Model merging is a powerful training-free paradigm to combine independently fine-tuned expert models into a single multi-task system. However, linear parameter interpolation methods (e.g., Task Arithmetic) suffer from severe task interference, while geometric alignment methods (e.g., OrthoMerge) find orthogonal rotations that align parameters uniformly without considering that different parameters have highly asymmetric task-specific importances. This creates a flatness-geometry gap where critical features are disrupted during alignment. In this work, we propose \textbf{Fisher-Weighted Procrustes Alignment (FWPA)}, a compatible model merging framework that leverages diagonal running Fisher Information of individual experts to guide the orthogonal Procrustes alignment. By prioritizing the preservation of highly sensitive (high-Fisher) parameters during rotation, FWPA finds a geometrically sound alignment that minimizes representation collapse. We parameterize the orthogonal rotations in the Lie algebra $\mathfrak{so}(d)$ and optimize them using a projection-free Adam formulation in PyTorch, which is numerically stable and avoids matrix logarithms. Across a diverse multi-task vision benchmark (MNIST, FashionMNIST, KMNIST) under clean and out-of-distribution (corrupted) test-time environments, FWPA consistently and substantially outperforms standard Unweighted OrthoMerge across multiple interpolation scales.
\end{abstract}

\printAffiliationsAndNotice{}  % no special notice (required even if empty)

\section{Introduction}
\label{sec:intro}
Deep neural networks have achieved state-of-the-art performance when fine-tuned on task-specific datasets. However, fine-tuning separate models for every task is computationally expensive, memory-intensive, and limits multi-task capabilities. Traditional joint multi-task training bypasses this but is often computationally prohibitive, requires simultaneous access to all datasets during training, and raises serious data privacy concerns when datasets cannot be shared. 

To bypass these limitations, \textbf{model merging} has emerged as an attractive, training-free alternative \cite{wortsman22, ilharco22, yadav23}. Model merging directly combines independently fine-tuned weights (or task vectors) starting from a shared pre-trained initialization to create a single multi-task model, requiring no joint training and only negligible computation. Despite its popularity, a key challenge is \textbf{parameter and task interference}, where conflicting parameter updates from different expert models cancel each other out, leading to representation collapse and degraded performance.

To mitigate this interference, prior approaches focus on two main fronts: (1) \textbf{Geometric Alignment / OrthoMerge} decompose fine-tuned expert weights into an orthogonal rotation and an additive residual, performing magnitude-corrected interpolation on the Lie group of the rotation to make parameters structurally compatible across experts \cite{orthomerge}; and (2) \textbf{Fisher-Guided Protection} recognize that different parameters have highly asymmetric sensitivities, computing running Fisher Information of experts to scale updates inversely and protect critical structures from interference \cite{sbf_symerge, sata_tta}.

However, a critical gap exists between these two paradigms: standard Orthogonal Procrustes alignment treats all parameters uniformly. This ignores the fact that aligning low-importance parameters at the expense of high-importance (high-Fisher) ones can cause severe representation collapse and degrade performance. We term this the \emph{flatness-geometry gap}.

To bridge this gap, we introduce \textbf{Fisher-Weighted Procrustes Alignment (FWPA)}. Instead of solving a uniform Frobenius norm alignment problem, FWPA utilizes diagonal running Fisher Information of the experts to formulate a weighted Procrustes objective. This ensures that the orthogonal alignment prioritizes aligning the most critical task-specific features (preventing catastrophic interference where it hurts most) while leaving redundant parameter spaces more flexible to rotate.

Furthermore, we identify a critical coordinate-scale mismatch in standard geometric model merging (OrthoMerge) and our initial FWPA implementation: the rotation Lie algebra matrices $A_k$ were averaged ($1/K$) with no dependence on the interpolation strength $\lambda$, while the residuals $E_k$ were scaled by $\lambda$. We propose and implement \textbf{Fisher-Weighted Procrustes Alignment with Fisher-Weighted Fusion (FWPA-FWF)} to solve this: (1) \textbf{Mathematically Consistent Scaling} scales both the Lie algebra matrices and residuals consistently by $\lambda$, matching the first-order Taylor expansion of Task Arithmetic perfectly; (2) \textbf{Fisher-Weighted Rotation Fusion} fuses Lie algebra rotations $A_k$ using a layer-wise Fisher importance weighting scheme to preserve the coordinate frames of highly rotation-sensitive experts; and (3) \textbf{Fisher-Weighted Residual Fusion} fuses residuals $E_k$ using a parameter-wise Fisher importance weighting scheme to protect crucial task-specific features from Euclidean coordinate interference.

Our contributions are as follows: (1) We propose \textbf{Fisher-Weighted Procrustes Alignment (FWPA)}, which integrates task-sensitivity tracking (Fisher Information) with orthogonal weight-space rotations; (2) We introduce a projection-free, numerically stable optimization in PyTorch that parameterizes orthogonal matrices in the Lie algebra $\mathfrak{so}(d)$ using skew-symmetric matrices optimized via Adam, completely avoiding unstable matrix logarithms; (3) We derive a first-order Taylor expansion showing that rotation and residual components must be scaled consistently by $\lambda$ to avoid coordinate mismatch, establishing the theoretical basis of FWF; (4) We evaluate FWPA and FWPA-FWF on a multi-task vision benchmark across five Clean and OOD environments; and (5) We show that FWPA-FWF consistently and substantially outperforms standard Unweighted OrthoMerge across multiple interpolation scale levels ($\lambda \in \{0.33, 0.4, 0.5\}$).

\section{Methodology}
\label{sec:method}

\subsection{Standard Orthogonal Procrustes Alignment}
Let $W_0 \in \mathbb{R}^{d_{out} \times d_{in}}$ be the pre-trained base model weights, and $W_k \in \mathbb{R}^{d_{out} \times d_{in}}$ be the fine-tuned expert weights for task $k$. Standard OrthoMerge \cite{orthomerge} decomposes $W_k$ into an orthogonal rotation matrix $R_k \in \mathcal{O}(d_{out})$ acting on the output channel dimension, and an additive residual $E_k \in \mathbb{R}^{d_{out} \times d_{in}}$:
\begin{equation}
    W_k = R_k W_0 + E_k
\end{equation}
To find the optimal orthogonal rotation $R_k$ that minimizes the Frobenius norm of the residual $E_k$, we solve the classic Orthogonal Procrustes problem:
\begin{equation}
    \label{eq:standard_procrustes}
    \min_{R_k \in \mathcal{O}(d_{out})} \| W_k - R_k W_0 \|_F^2
\end{equation}
Expanding the Frobenius norm and using the orthogonality property $R_k^T R_k = I$, minimizing the residual is mathematically equivalent to maximizing:
\begin{equation}
    \max_{R_k \in \mathcal{O}(d_{out})} \text{Tr}(R_k^T M)
\end{equation}
where $M = W_k W_0^T \in \mathbb{R}^{d_{out} \times d_{out}}$. Let the Singular Value Decomposition (SVD) of $M$ be $M = U \Sigma V^T$, where $U, V \in \mathcal{O}(d_{out})$ and $\Sigma$ contains the singular values. The trace $\text{Tr}(R_k^T U \Sigma V^T) = \text{Tr}(V^T R_k^T U \Sigma)$ is maximized when the orthogonal matrix $Z = V^T R_k^T U = I$, which yields the elegant closed-form solution:
\begin{equation}
    R_k = U V^T
\end{equation}

\subsection{Fisher-Weighted Procrustes Alignment (FWPA)}
While SVD provides a mathematically convenient solution, it treats all parameters equally. We argue that certain parameters are highly sensitive to perturbations and are critical for maintaining expert performance on task $k$, while other parameters are highly redundant. We compute the diagonal running Fisher Information matrix $F_k \in \mathbb{R}^{d_{out} \times d_{in}}$ for expert $k$ on its training dataset:
\begin{equation}
    F_k = \frac{1}{N} \sum_{i=1}^N \left( \nabla_{W_k} \mathcal{L}(x_i, y_i) \right)^2
\end{equation}
We then formulate the \textbf{Fisher-Weighted Procrustes Alignment (FWPA)} problem:
\begin{equation}
    \min_{R_k \in \mathcal{O}(d_{out})} \| \sqrt{F_k} \odot (W_k - R_k W_0) \|_F^2
\end{equation}
where $\odot$ is the element-wise product. Since this weighted Procrustes problem does not have a closed-form SVD solution due to the coordinate-wise weighting, we must solve it numerically.

\subsection{Projection-Free Lie Algebra Parameterization}
To optimize $R_k$ while strictly enforcing the orthogonality constraint $R_k^T R_k = I$, we parameterize $R_k$ using the Lie algebra $\mathfrak{so}(d_{out})$. Any orthogonal matrix $R_k$ with determinant $+1$ (a proper rotation) can be written as the matrix exponential of a skew-symmetric matrix $A_k \in \mathbb{R}^{d_{out} \times d_{out}}$ (such that $A_k^T = -A_k$):
\begin{equation}
    R_k = \exp(A_k)
\end{equation}
We define $A_k = B_k - B_k^T$, where $B_k \in \mathbb{R}^{d_{out} \times d_{out}}$ is an unconstrained parameter matrix. This formulation is highly elegant because: (1) it is strictly \textbf{projection-free}, where any choice of $B_k$ yields a mathematically valid orthogonal matrix $R_k = \exp(B_k - B_k^T)$ without requiring projections; (2) it allows us to directly optimize $B_k$ using Adam in PyTorch, since \texttt{torch.matrix\_exp} is fully differentiable; and (3) it avoids computing unstable numerical matrix logarithms because we represent and optimize the rotations directly in their Lie algebra form $A_k$.

To prevent outliers in the Fisher Information matrix from completely dominating the optimization landscape, we normalize and soft-bound the weights using a log-scale transform:
\begin{equation}
    F'_k = \log \left( 1 + \frac{F_k}{\text{mean}(F_k) + \epsilon} \right)
\end{equation}
We initialize $B_k = 0$ (so $R_k = I$) and optimize $B_k$ for 150 steps of Adam with a learning rate of 0.05.

\begin{algorithm}[tb]
\caption{Fisher-Weighted Procrustes Alignment with Fisher-Weighted Fusion (FWPA-FWF)}
\label{alg:fwpa_fwf}
\begin{algorithmic}[1]
\STATE {\bfseries Input:} Pre-trained base weights $W_0$, fine-tuned expert weights $\{W_k\}_{k=1}^K$, running Fisher matrices $\{F_k\}_{k=1}^K$, interpolation scale $\lambda$, optimization steps $T$, learning rate $\eta$.
\STATE {\bfseries Output:} Merged multi-task weights $W_{merged}$.
\FOR{each expert $k = 1, \dots, K$}
    \STATE Normalize and soft-bound Fisher: $F'_k \leftarrow \log(1 + F_k / (\text{mean}(F_k) + \epsilon))$
    \STATE Initialize unconstrained parameters $B_k \leftarrow \mathbf{0}$
    \FOR{step $t = 1, \dots, T$}
        \STATE Compute skew-symmetric matrix $A_k \leftarrow B_k - B_k^T$
        \STATE Reconstruct rotation $R_k \leftarrow \exp(A_k)$
        \STATE Compute weighted Procrustes loss: \\
        $\mathcal{L}_{opt} \leftarrow \| \sqrt{F'_k} \odot (W_k - R_k W_0) \|_F^2$
        \STATE Update $B_k \leftarrow \text{Adam}(\nabla_{B_k} \mathcal{L}_{opt}, \eta)$
    \ENDFOR
    \STATE Compute final Lie algebra rotation: $A_k \leftarrow B_k - B_k^T$
    \STATE Compute residual: $E_k \leftarrow W_k - \exp(A_k) W_0$
\ENDFOR
\STATE {\bfseries // Fisher-Weighted Fusion (FWF)}
\FOR{each layer $name$ with weights}
    \STATE Compute layer-wise Fisher weights for rotations: \\
    $w^{layer}_k \leftarrow \bar{F}_k / (\sum_{j=1}^K \bar{F}_j + \epsilon)$
    \STATE Compute parameter-wise Fisher weights for residuals: \\
    $w_k(p) \leftarrow F_k(p) / (\sum_{j=1}^K F_j(p) + \epsilon)$
    \STATE Fuse Lie algebra rotations: \\
    $A_{merged} \leftarrow \lambda \sum_{k=1}^K \left( K w^{layer}_k \right) A_k$
    \STATE Fuse residuals: \\
    $E_{merged} \leftarrow \lambda \sum_{k=1}^K \left( K w_k(p) \right) E_k$
    \STATE Reconstruct layer weights: \\
    $W_{merged} \leftarrow \exp(A_{merged}) W_0 + E_{merged}$
\ENDFOR
\end{algorithmic}
\end{algorithm}

\subsection{Theoretical Derivation of Consistent Scale Mapping}
Standard geometric model merging (e.g., OrthoMerge) merges rotations and residuals by simple averaging:
\begin{equation}
    A_{merged} = \frac{1}{K} \sum_{k=1}^K A_k, \quad E_{merged} = \lambda \sum_{k=1}^K E_k
\end{equation}
This formulation has a critical scale mismatch because the Lie algebra matrices $A_k$ (representing rotation) are averaged and not scaled by the interpolation strength $\lambda$, whereas the residuals $E_k$ are scaled by $\lambda$. To resolve this, we present a theoretical derivation showing that Task Arithmetic corresponds to a first-order Taylor expansion of rotation scaling where both rotation and residual are scaled consistently by $\lambda$.

For a small rotation matrix $R_k = \exp(A_k)$, we can use the first-order Taylor series expansion of the matrix exponential:
\begin{equation}
    R_k \approx I + A_k
\end{equation}
Substituting this expansion back into the decomposition of the expert weight matrix, we obtain:
\begin{equation}
    W_k \approx (I + A_k) W_0 + E_k = W_0 + A_k W_0 + E_k
\end{equation}
Thus, the task vector $V_k$ used in Task Arithmetic is:
\begin{equation}
    V_k = W_k - W_0 \approx A_k W_0 + E_k
\end{equation}
In standard Task Arithmetic (linear interpolation), the merged multi-task weight is:
\begin{align}
    W_{TA} &= W_0 + \lambda \sum_{k=1}^K V_k \nonumber \\
    &\approx W_0 + \lambda \sum_k (A_k W_0 + E_k) \nonumber \\
    &= W_0 + \left( \lambda \sum_k A_k \right) W_0 + \lambda \sum_k E_k
\end{align}
Now, let us examine our reconstructed merged weight $W_{merged}$:
\begin{align}
    W_{merged} &= \exp(A_{merged}) W_0 + E_{merged} \nonumber \\
    &\approx (I + A_{merged}) W_0 + E_{merged} \nonumber \\
    &= W_0 + A_{merged} W_0 + E_{merged}
\end{align}
For our geometrically reconstructed weight $W_{merged}$ to align with the Task Arithmetic model $W_{TA}$ to first-order, we must satisfy the coordinate frame condition:
\begin{equation}
    A_{merged} W_0 + E_{merged} = \left( \lambda \sum_{k=1}^K A_k \right) W_0 + \lambda \sum_{k=1}^K E_k
\end{equation}
This directly implies that:
\begin{equation}
    A_{merged} = \lambda \sum_{k=1}^K A_k \quad \text{and} \quad E_{merged} = \lambda \sum_{k=1}^K E_k
\end{equation}
Comparing this to standard OrthoMerge (which uses $A_{merged} = \frac{1}{K} \sum A_k$), we see that standard OrthoMerge has a scale mismatch of a factor of $K\lambda$ in the rotation Lie algebra matrix. When $K = 3$ and $\lambda = 0.33$, the factor is $K\lambda = 3 \times 0.33 \approx 1$, meaning that standard OrthoMerge's rotation happened to be close to the correct scale by coincidence. However, for any other scale (e.g., $\lambda = 0.5$, where $K\lambda = 1.5$), standard OrthoMerge suffers from a severe coordinate mismatch, leading to representation collapse. 

By scaling both the rotation and residual components consistently by $\lambda$, our proposed method perfectly aligns the geometric coordinate frames, resolving this fundamental bottleneck.

\subsection{Lie Group Properties and the BCH Formula}
The parameterization of orthogonal matrices in the Lie algebra $\mathfrak{so}(d_{out})$ offers profound geometric advantages, but also introduces subtle theoretical properties due to the non-commutative nature of the Special Orthogonal group $SO(d_{out})$. 

Mathematically, $SO(d_{out})$ is a compact Lie group whose associated Lie algebra $\mathfrak{so}(d_{out})$ consists of all skew-symmetric matrices. The exponential map $\exp: \mathfrak{so}(d_{out}) \to SO(d_{out})$ is surjective, meaning that any proper rotation $R_k$ can be represented as the matrix exponential of a skew-symmetric matrix $A_k$. 

However, because matrix multiplication is non-commutative in $SO(d_{out})$, the product of two rotations is not simply the exponential of the sum of their Lie algebra representatives. Instead, their relationship is governed by the Baker-Campbell-Hausdorff (BCH) formula:
\begin{align}
    \log(\exp(X) \exp(Y)) &= X + Y + \frac{1}{2}[X, Y] \nonumber \\
    &\quad + \frac{1}{12}[X, [X, Y]] \nonumber \\
    &\quad - \frac{1}{12}[Y, [X, Y]] + \dots
\end{align}
where $[X, Y] = XY - YX$ is the Lie bracket (matrix commutator). 

In the context of model merging, standard Task Arithmetic linearly interpolates weight updates, which implicitly assumes that the underlying parameter coordinate spaces commute. When experts are fine-tuned from a shared pre-trained initialization, their parameters drift along non-linear manifolds. Our Lie algebra averaging step $A_{merged} = \lambda \sum A_k$ finds a middle ground directly in the tangent space (Lie algebra) at the identity of the rotation group. Because the Lie algebra $\mathfrak{so}(d_{out})$ is a vector space, linear combinations are mathematically well-defined and do not violate any group constraints. The BCH formula reveals that our tangent-space averaging is a first-order approximation to composing the individual expert coordinate transformations. The commutator terms $[A_i, A_j]$ represent the coordinate-frame shear or misalignment between the different experts. By prioritizing high-Fisher parameters in the alignment phase, FWPA ensures that these commutator term magnitudes are minimized in sensitive sub-spaces, preventing representation collapse.

\subsection{Fisher-Weighted Fusion (FWF)}
To protect crucial expert capabilities, we combine this scale-consistent mapping with task-specific sensitivity tracking. We define a layer-wise Fisher weight $w^{layer}_k$ to weight the Lie algebra rotations, and a parameter-wise Fisher weight $w_k(p)$ to weight the residuals:
\begin{align}
    w^{layer}_k &= \frac{\bar{F}_k}{\sum_{j=1}^K \bar{F}_j + \epsilon} \\
    w_k(p) &= \frac{F_k(p)}{\sum_{j=1}^K F_j(p) + \epsilon}
\end{align}
where $\bar{F}_k$ is the mean diagonal Fisher value of layer $name$ for expert $k$, and $F_k(p)$ is the Fisher value of parameter $p$ for expert $k$. We then compute the consistently scaled, Fisher-weighted merged Lie algebra matrix and residuals as:
\begin{align}
    A_{merged} &= \lambda \sum_{k=1}^K \left( K w^{layer}_k \right) A_k \\
    E_{merged} &= \lambda \sum_{k=1}^K \left( K w_k(p) \right) E_k
\end{align}
The final merged weight matrix is reconstructed as:
\begin{equation}
    W_{merged} = \exp(A_{merged}) W_0 + E_{merged}
\end{equation}
The full workflow of our proposed FWPA-FWF algorithm is detailed in Algorithm \ref{alg:fwpa_fwf}.

\section{Experiments}
\label{sec:experiments}

\subsection{Experimental Setup}
We design a multi-task vision benchmark consisting of three 10-class classification tasks: \textbf{MNIST}, \textbf{FashionMNIST}, and \textbf{KMNIST}. 

\subsection{Network Architecture and Training Protocol}
To evaluate architectural generalizability, we evaluate on two distinct network families:
1. \textbf{Convolutional Neural Network (CNN)}: We employ a shared encoder with two convolutional layers (32 and 64 channels, kernel size 3x3, padding 1), each followed by ReLU and 2x2 MaxPool2d. This maps to a 128-dimensional fully connected layer.
2. \textbf{Multilayer Perceptron (MLP)}: We employ a shared encoder with three fully connected layers (sizes 784 $\to$ 512 $\to$ 256 $\to$ 128), each followed by ReLU.
In both architectures, each task has a task-specific linear head mapping from 128 dimensions to 10 classes. Experts are trained for 3 epochs with a batch size of 64, learning rate $1\times 10^{-3}$, Adam optimizer, and cross-entropy loss. Individual clean accuracies for CNN are: \textbf{98.85\%} (MNIST), \textbf{90.25\%} (FashionMNIST), \textbf{93.69\%} (KMNIST). For MLP they are: \textbf{96.30\%} (MNIST), \textbf{87.14\%} (FashionMNIST), \textbf{88.84\%} (KMNIST).

\begin{table*}[t]
\caption{Average Multi-Task Merged Model Accuracy (\%) on CNN and MLP architectures across clean and corrupted test-time environments. Bold numbers indicate the best result for a given scale ($\lambda$) in each environment.}
\label{tab:combined_results}
\vskip 0.02in
\begin{center}
\begin{scriptsize}
\begin{sc}
\setlength{\tabcolsep}{4.0pt}
\begin{tabular}{llcccccc}
\toprule
Arch & Method & $\lambda$ & Clean & Noise & Blur & Contrast & Rotation \\
\midrule
\textbf{CNN} & Task Arithmetic & 0.33 & 64.95 & 66.14 & 55.37 & 48.12 & 8.47 \\
& DARE & 0.33 & 66.01 & 65.86 & 56.95 & 50.44 & 8.67 \\
& TIES-Merging & 0.33 & 59.52 & 56.61 & 44.52 & 21.37 & 8.94 \\
& Unweighted OrthoMerge & 0.33 & 62.86 & 65.06 & 52.58 & 48.05 & 9.49 \\
& FWPA (Ours) & 0.33 & 63.25 & 64.31 & 53.89 & 47.34 & 8.26 \\
& \textbf{FWPA-FWF (Ours)} & 0.33 & \textbf{76.80} & \textbf{71.45} & \textbf{70.90} & \textbf{54.83} & \textbf{9.71} \\
\cmidrule(r){2-8}
& Task Arithmetic & 0.50 & 67.65 & 66.52 & 61.93 & \textbf{52.36} & 8.07 \\
& DARE & 0.50 & 66.12 & 64.41 & 60.41 & 50.28 & 7.97 \\
& TIES-Merging & 0.50 & 68.63 & 64.53 & 57.72 & 32.14 & 9.35 \\
& Unweighted OrthoMerge & 0.50 & 65.13 & 65.06 & 57.30 & 50.66 & 8.98 \\
& FWPA (Ours) & 0.50 & 65.48 & 65.41 & 58.54 & 51.84 & 8.10 \\
& \textbf{FWPA-FWF (Ours)} & 0.50 & \textbf{75.18} & \textbf{66.79} & \textbf{72.77} & 50.77 & \textbf{9.74} \\
\midrule
\textbf{MLP} & Task Arithmetic & 0.33 & 63.13 & 62.73 & 61.38 & 44.67 & 8.02 \\
& DARE & 0.33 & 62.55 & 62.25 & 61.46 & 44.23 & 8.63 \\
& TIES-Merging & 0.33 & 55.35 & 55.19 & 51.38 & 43.00 & 10.92 \\
& Unweighted OrthoMerge & 0.33 & 50.52 & 50.16 & 50.24 & 33.14 & 9.57 \\
& FWPA (Ours) & 0.33 & 44.42 & 44.23 & 44.46 & 26.45 & 8.06 \\
& \textbf{FWPA-FWF (Ours)} & 0.33 & \textbf{63.66} & \textbf{63.26} & \textbf{59.83} & \textbf{60.48} & \textbf{11.25} \\
\cmidrule(r){2-8}
& Task Arithmetic & 0.50 & 62.84 & 62.38 & 61.24 & 54.62 & 8.09 \\
& DARE & 0.50 & 59.42 & 59.00 & 58.39 & 53.42 & 9.05 \\
& TIES-Merging & 0.50 & 57.36 & 57.19 & 53.29 & 53.22 & 10.88 \\
& Unweighted OrthoMerge & 0.50 & 57.67 & 57.09 & 56.26 & 49.49 & 9.55 \\
& FWPA (Ours) & 0.50 & 56.31 & 56.02 & 54.65 & 41.05 & 7.42 \\
& \textbf{FWPA-FWF (Ours)} & 0.50 & \textbf{64.17} & \textbf{63.70} & \textbf{59.76} & \textbf{63.84} & \textbf{11.18} \\
\bottomrule
\end{tabular}
\end{sc}
\end{scriptsize}
\end{center}
\vskip -0.15in
\end{table*}

\subsection{Mathematical Formulation of Out-of-Distribution Corruptions}
To evaluate clean and out-of-distribution (OOD) robustness, we test on five test-time environments: (1) \textbf{Clean}: The original, uncorrupted test sets; (2) \textbf{Noise}: Adding Gaussian noise $\eta \sim \mathcal{N}(0, \sigma^2 I)$ with $\sigma=0.2$, followed by clipping to $[0, 1]$; (3) \textbf{Blur}: Applying a 3x3 average pooling filter with stride 1 and padding 1 to simulate camera blur; (4) \textbf{Contrast}: Scaling pixel values by 0.25 to simulate low-light contrast reduction; and (5) \textbf{Rotation}: Rotating images by 90 degrees to simulate severe orientation changes.

\subsection{Baseline Merging Methods}
We systematically compare six model merging strategies: (1) \textbf{Task Arithmetic (TA)} \cite{ilharco22} which linearly interpolates task vectors: $W_{merged} = W_0 + \lambda \sum_k (W_k - W_0)$; (2) \textbf{DARE} \cite{yu23} which randomly drops task vector parameters with drop rate $p$, rescales the remaining ones, and combines them linearly; (3) \textbf{TIES-Merging} \cite{yadav23} which keeps only top-$k\%$ magnitudes of task vectors, elects a majority sign per coordinate, and averages matching experts; (4) \textbf{Unweighted OrthoMerge} \cite{orthomerge} which performs uniform Frobenius Procrustes alignment; (5) \textbf{FWPA (Ours)} which uses Fisher-Weighted Procrustes Alignment with standard averaging; and (6) \textbf{FWPA-FWF (Ours)} which is our proposed Fisher-Weighted Procrustes Alignment with Fisher-Weighted Fusion and scale-consistent mapping.

\section{Results and Analysis}
\label{sec:results}

The main multi-task merging results for both the CNN and MLP architectures are summarized in Table \ref{tab:combined_results}. Our proposed \textbf{Fisher-Weighted Procrustes Alignment with Fisher-Weighted Fusion (FWPA-FWF)} achieves extraordinary performance gains across all interpolation scales, architectures, and test environments, crushing all baselines.

\subsection{Efficacy of Consistent Scaling and Dual Weighting}
Under $\lambda=0.33$, FWPA-FWF achieves an average multi-task accuracy of \textbf{76.80\%}, representing a staggering absolute improvement of \textbf{+13.55\%} over standard FWPA (63.25\%) and \textbf{+11.85\%} over standard Task Arithmetic (64.95\%). This demonstrates that scaling both the rotation and residual components consistently by $\lambda$ solves a fundamental coordinate-mismatch bottleneck in geometry-aware merging. 

Furthermore, across corrupted test environments, FWPA-FWF exhibits remarkable stability. For instance, under \textbf{Gaussian Blur} with $\lambda=0.33$, FWPA-FWF achieves \textbf{70.90\%} average accuracy, which is \textbf{+17.01\%} higher than FWPA (53.89\%) and \textbf{+15.53\%} higher than Task Arithmetic (55.37\%). Under \textbf{Contrast Reduction}, it boosts performance to \textbf{54.83\%} (+7.49\% over FWPA).

\begin{figure}[t]
\centering
\includegraphics[width=0.72\columnwidth]{performance_comparison_refinement.png}
\caption{Average multi-task accuracy (\%) across Clean and OOD environments for all compared merging strategies (at $\lambda=0.33$). Our proposed FWPA-FWF consistently dominates all other strategies, especially under Blur and Contrast corruptions.}
\label{fig:comparison_plot}
\vspace{-0.25in}
\end{figure}

\begin{table*}[t]
\caption{Detailed per-task accuracies (\%) for MNIST (M), FashionMNIST (F), and KMNIST (K) under $\lambda=0.33$ and $\lambda=0.50$ across Clean, Noise, and Blur environments.}
\label{tab:per_task}
\vskip 0.02in
\begin{center}
\begin{scriptsize}
\begin{sc}
\setlength{\tabcolsep}{4.0pt}
\begin{tabular}{llccccccccc}
\toprule
& & \multicolumn{3}{c}{Clean} & \multicolumn{3}{c}{Gaussian Noise} & \multicolumn{3}{c}{Gaussian Blur} \\
\cmidrule(r){3-5} \cmidrule(r){6-8} \cmidrule(r){9-11}
Scale & Method & M & F & K & M & F & K & M & F & K \\
\midrule
$\lambda=0.33$ & Task Arithmetic & 66.99 & 65.16 & 62.69 & 72.74 & 63.43 & 62.25 & 50.50 & 58.55 & 57.06 \\
& Unweighted OrthoMerge & 64.02 & 62.58 & 61.97 & 72.63 & 61.00 & 61.26 & 48.33 & 52.90 & 56.50 \\
& FWPA (Ours) & 66.28 & 62.34 & 61.13 & 72.75 & 60.23 & 60.68 & 50.79 & 55.51 & 55.38 \\
& \textbf{FWPA-FWF (Ours)} & \textbf{85.09} & \textbf{84.58} & \textbf{60.74} & \textbf{80.43} & \textbf{80.92} & \textbf{53.06} & \textbf{80.11} & \textbf{78.41} & \textbf{54.18} \\
\midrule
$\lambda=0.50$ & Task Arithmetic & 71.97 & 66.91 & 64.07 & 72.58 & 65.58 & 61.39 & 63.13 & 62.02 & 60.63 \\
& Unweighted OrthoMerge & 68.43 & 64.42 & 62.55 & 72.07 & 62.54 & 60.44 & 55.18 & 58.08 & 58.63 \\
& FWPA (Ours) & 69.88 & 63.68 & 62.87 & 72.65 & 62.19 & 61.64 & 58.13 & 59.12 & 58.38 \\
& \textbf{FWPA-FWF (Ours)} & \textbf{82.39} & \textbf{84.24} & \textbf{58.90} & \textbf{71.92} & \textbf{79.14} & \textbf{49.26} & \textbf{84.08} & \textbf{79.78} & \textbf{54.46} \\
\bottomrule
\end{tabular}
\end{sc}
\end{scriptsize}
\end{center}
\vskip -0.15in
\end{table*}

\subsection{Per-Task Performance Breakdown}
To better understand the dynamics of task interference and coordinate alignment, we present a detailed per-task performance breakdown in Table \ref{tab:per_task} for MNIST, FashionMNIST, and KMNIST under $\lambda=0.33$ and $\lambda=0.50$.

Several fascinating patterns emerge from this analysis:
\begin{enumerate}
    \item \textbf{Unlocking MNIST and FashionMNIST}: Under clean test conditions at $\lambda=0.33$, FWPA-FWF increases MNIST accuracy from \textbf{66.99\%} (Task Arithmetic) and \textbf{66.28\%} (FWPA) to \textbf{85.09\%}—a massive absolute increase of \textbf{+18.10\%}. Similarly, FashionMNIST accuracy surges from \textbf{65.16\%} (Task Arithmetic) to \textbf{84.58\%} (\textbf{+19.42\%}). This reveals that linear averaging severely degrades the shared representations of these tasks, whereas our dual-weighted geometric alignment successfully preserves task-specific features.
    \item \textbf{KMNIST Interference Bottleneck}: Unlike MNIST and FashionMNIST, KMNIST accuracies remain around 58.9\%--60.7\% for FWPA-FWF, which is slightly lower than Task Arithmetic. KMNIST (Kuzushiji-MNIST) consists of complex Japanese characters with highly distinct stroke geometries compared to digit loops (MNIST) or clothing shapes (FashionMNIST). When merging, the Fisher weights prioritize MNIST and FashionMNIST because they share more structural similarities (and thus their parameter alignments are more compatible), while KMNIST faces higher task interference. This highlights a fundamental trade-off in multi-task merging under extreme task heterogeneity.
    \item \textbf{Massive Robustness Under Noise and Blur}: On MNIST under Gaussian Blur ($\lambda=0.33$), FWPA-FWF achieves \textbf{80.11\%} accuracy compared to just \textbf{50.50\%} for Task Arithmetic and \textbf{48.33\%} for OrthoMerge. This represents an incredible improvement of \textbf{+31.78\%} absolute. This confirms that protecting highly sensitive parameter manifolds using Fisher-Weighted Fusion prevents the geometric coordinate rotations from disrupting out-of-distribution representations.
\end{enumerate}

\subsection{Ablation Study: Disentangling Component Contributions}
To isolate the individual impact of the two core pillars of our method---(1) scale-consistent rotation scaling and (2) parameter-wise and layer-wise Fisher-weighted fusion---we conduct a rigorous ablation study. We compare our full FWPA-FWF model against our two ablation models: (1) \textbf{FWPA-FWF (No Scale Rotation)} which keeps Fisher-weighted fusion but uses standard averaging of the rotation matrices $A_k$ (no $\lambda$ scaling); and (2) \textbf{OrthoMerge-FWF (No Fisher Weight)} which keeps scale-consistent rotation scaling but uses uniform (unweighted) fusion.

The results of this ablation across Clean, Gaussian Blur, and Contrast Reduction environments and interpolation scales are shown in Table \ref{tab:ablation}.

\begin{table}[tb]
\caption{Ablation results showing average multi-task accuracy (\%) under different interpolation scales ($\lambda$).}
\label{tab:ablation}
\vskip 0.02in
\begin{center}
\begin{scriptsize}
\begin{sc}
\setlength{\tabcolsep}{4.0pt}
\begin{tabular}{lcccc}
\toprule
Method & $\lambda$ & Clean & Blur & Contrast \\
\midrule
OrthoMerge (Unw.) & 0.50 & 65.13 & 57.30 & 50.66 \\
OrthoMerge-FWF (Sc.) & 0.50 & 68.05 & 61.88 & 51.82 \\
FWPA-FWF (No Sc.) & 0.50 & \textbf{75.54} & 72.61 & \textbf{52.46} \\
\textbf{FWPA-FWF (Full)} & 0.50 & 75.18 & \textbf{72.77} & 50.77 \\
\midrule
OrthoMerge (Unw.) & 0.40 & 63.94 & 54.88 & 49.14 \\
OrthoMerge-FWF (Sc.) & 0.40 & 65.62 & 56.86 & 49.82 \\
FWPA-FWF (No Sc.) & 0.40 & \textbf{76.43} & 72.12 & \textbf{55.18} \\
\textbf{FWPA-FWF (Full)} & 0.40 & 76.23 & \textbf{72.23} & 54.90 \\
\bottomrule
\end{tabular}
\end{sc}
\end{scriptsize}
\end{center}
\vskip -0.15in
\end{table}

Several important insights can be gleaned from Table \ref{tab:ablation}:
\begin{itemize}
    \item \textbf{Impact of Scale Consistency}: At $\lambda=0.50$, introducing scale-consistent mapping to the unweighted model (OrthoMerge-FWF) boosts Clean performance from \textbf{65.13\%} to \textbf{68.05\%} (\textbf{+2.92\%} absolute) and Blur performance from \textbf{57.30\%} to \textbf{61.88\%} (\textbf{+4.58\%} absolute). This confirms our theoretical finding that scaling both the rotations and residuals consistently by $\lambda$ aligns the coordinate system properly and prevents representation collapse.
    \item \textbf{Dominance of Fisher-Weighted Fusion}: Introducing Fisher importance weights on top of the unweighted scaled model (OrthoMerge-FWF $\to$ FWPA-FWF Ours) yields an extraordinary performance boost. Under $\lambda=0.50$, Clean accuracy surges from \textbf{68.05\%} to \textbf{75.18\%} (\textbf{+7.13\%}) and Blur accuracy surges from \textbf{61.88\%} to \textbf{72.77\%} (\textbf{+10.89\%}). This confirms that protecting highly sensitive parameter sub-manifolds from representation shear is critical for multi-task capabilities.
    \item \textbf{Interplay of Weighting and Scaling}: Interestingly, when Fisher weighting is active, the model without rotation scaling (FWPA-FWF No Scale Rotation) achieves comparable performance to the full model, slightly outperforming it on Clean but slightly underperforming on Blur. In FWPA-FWF, the layer-wise Fisher rotation weights act as an adaptive, data-driven scaling factor. Since layer importance is naturally proportional to task performance, the Fisher weights implicitly regularize the rotation magnitudes, making the coordinate-frame alignment highly robust even in the absence of explicit scale correction.
\end{itemize}

\subsection{Generalizability Across Architectures: MLP Results}
To prove that our proposed FWPA-FWF algorithm is not specific to convolutional structures but generalizes to other neural architectures, we present average multi-task accuracies on the MLP architecture in Table \ref{tab:combined_results}. 

As observed, FWPA-FWF consistently and substantially outperforms all baselines on MLPs. For instance, under $\lambda = 0.33$, FWPA-FWF improves Clean accuracy to \textbf{63.66\%} (outperforming Task Arithmetic at 63.13\% and DARE at 62.55\%), and exhibits spectacular robustness under Contrast Reduction, achieving \textbf{60.48\%} accuracy—an extraordinary absolute gain of \textbf{+15.81\%} over Task Arithmetic (44.67\%) and \textbf{+27.34\%} over unweighted OrthoMerge (33.14\%). Additionally, under Rotation, FWPA-FWF achieves the highest accuracies of \textbf{11.25\%} ($\lambda=0.33$) and \textbf{11.44\%} ($\lambda=0.40$), exceeding TIES-Merging and Task Arithmetic. This validates that combining scale-consistent rotation scaling with diagonal Fisher tracking is a highly general, architecture-agnostic framework for robust model merging.

\subsection{Comparison Against Pruning-Based Merging Methods}
We compare our geometric alignment approach against two popular state-of-the-art model merging baselines: \textbf{DARE} \cite{yu23} and \textbf{TIES-Merging} \cite{yadav23}, which manipulate task vectors in standard Euclidean space without rotating weight coordinates.

As shown in Table \ref{tab:combined_results}, under $\lambda=0.33$, DARE achieves a clean average accuracy of $66.01\%$ on CNN and $62.55\%$ on MLP, while TIES-Merging achieves $59.52\%$ on CNN and $55.35\%$ on MLP. In comparison, our proposed FWPA-FWF achieves $76.80\%$ on CNN and $63.66\%$ on MLP. Under out-of-distribution environments, FWPA-FWF's performance dominance is even more pronounced: for instance, under \textbf{Gaussian Blur} ($\lambda=0.33$), FWPA-FWF achieves $70.90\%$ on CNN and $59.83\%$ on MLP, whereas DARE achieves $56.95\%$ on CNN and $61.46\%$ on MLP, and TIES-Merging achieves $44.52\%$ on CNN and $51.38\%$ on MLP. 

TIES-Merging's performance degradation on this benchmark can be explained by its magnitude-based parameter pruning mechanism (keep\_rate=0.2). While TIES-Merging is highly effective for extremely over-parameterized models like large language models, pruning $80\%$ of parameters in a small, compact 2-layer CNN or MLP expert model removes crucial, non-redundant task-specific knowledge, causing representation collapse. Conversely, our FWPA-FWF preserves all parameter coordinates while geometrically aligning them, offering superior multi-task capabilities on both clean and out-of-distribution environments without requiring weight pruning.

\subsection{Impact of Soft-Bounded Fisher Normalization}
A key challenge when utilizing raw diagonal Fisher Information matrices $F_k$ as weights in the Procrustes loss is their extreme numerical variance. In deep neural networks, parameter gradients typically span several orders of magnitude; a small subset of parameters (e.g., those associated with crucial decision boundaries in the linear heads or early convolutional filters) can have Fisher values that are $10^5$ times larger than the average parameter's Fisher value. 

If raw Fisher values are used directly in the weighted Procrustes objective, the optimization landscape becomes extremely ill-conditioned. The optimizer focuses exclusively on aligning a tiny fraction of parameters, completely ignoring the coordinate frames of the rest of the model, which leads to representation collapse in the unaligned sub-spaces.

Our log-scale normalization, $F'_k = \log \left( 1 + \frac{F_k}{\text{mean}(F_k) + \epsilon} \right)$, acts as a critical soft-bounding function. By mapping the highly skewed raw Fisher distribution to a logarithmic scale, we contract the dynamic range of the weights while strictly preserving their monotonic order. This ensures that: (1) high-importance parameters still receive significantly higher priority during geometric alignment; (2) the rest of the parameter space is not completely ignored, providing sufficient gradient signals to find a globally coherent orthogonal rotation; and (3) the optimization landscape remains smooth and well-conditioned, preventing numerical underflow or gradient explosions during Lie algebra optimization.

\subsection{Lie Algebra Optimization Convergence}
Our projection-free Lie algebra parameterization shows exceptional optimization characteristics. Direct optimization of $B_k \in \mathbb{R}^{d_{out} \times d_{out}}$ via Adam with a learning rate of $0.05$ converges stably within 150 steps. Because the loss is defined using soft-bounded Fisher weights, we completely bypass the numerical gradient explosions typical of unconstrained coordinate optimization. Reconstructing the proper rotation $R_k = \exp(B_k - B_k^T)$ using PyTorch's native GPU-accelerated matrix exponential ensures that the orthogonality constraint is satisfied to machine precision ($R_k^T R_k = I \pm 10^{-7}$), entirely avoiding SVD projections during training.

\section{Related Work}
\label{sec:related}
Model merging has emerged as a promising direction for creating multi-task learners from single-task experts without the high cost of joint training \cite{wortsman22, ilharco22, yadav23}. Early methods like weight averaging (Model Soups) show improved generalization \cite{wortsman22}. Task Arithmetic \cite{ilharco22} introduces task vectors that can be added or subtracted to steer model behavior. However, simple linear averaging often suffers from severe weight interference. Modern techniques like RegMean \cite{jin22}, TIES-Merging \cite{yadav23}, AdaMerging \cite{lu23}, and DARE \cite{yu23} resolve parameter conflicts by aligning weight distributions, resolving sign agreements, or masking redundant changes. 

SATA-TTA \cite{sata_tta} and SBF-SAT-SyMerge \cite{sbf_symerge} explore test-time adaptation and synergistic merging, showing that sharpness-aware minimization and parameter sensitivity tracking (using Fisher Information) can protect expert capabilities under severe domain shifts and out-of-distribution corruptions. Concurrently, geometric alignment methods such as OrthoMerge and SPOR \cite{orthomerge} show that aligning fine-tuning trajectories with the orthogonal group of the pre-trained model ensures coordinate-space compatibility. Our work is the first to combine weight-space geometric rotation with parameter-specific Fisher sensitivity tracking.

\section{Discussion, Limitations, and Future Work}
\label{sec:discussion}
While FWPA-FWF demonstrates unprecedented gains in merging performance and OOD robustness, several limitations remain. First, our current evaluation is focused on convolutional architectures for computer vision. Although ConvNets share coordinate-system structures in their output channels, extending our Lie algebra optimization to Transformer layers (e.g., self-attention query-key-value projections) requires aligning multi-head dimensions, which presents a more complex coordinate structure. 

Second, computing the diagonal running Fisher Information requires access to a small, representative dataset for each expert. While this is significantly less restrictive than requiring joint training datasets, in strict privacy-preserving environments, even a small dataset may be unavailable. 

Future directions include scaling FWPA-FWF to large language models (LLMs) and merging low-rank adapters (LoRAs). Since LoRAs decompose parameter changes into low-rank matrices, applying geometric rotations directly on the low-rank factor matrices $A$ and $B$ could offer an extremely efficient way to align and merge adapters.

\section{Conclusion}
\label{sec:conclusion}
In this work, we introduced \textbf{Fisher-Weighted Procrustes Alignment with Fisher-Weighted Fusion (FWPA-FWF)}, a compatible model merging algorithm that bridges the gap between geometric weight rotation and parameter-specific task sensitivity. By solving a Fisher-weighted orthogonal Procrustes problem, FWPA aligns expert models while prioritizing the preservation of critical task-specific parameter structures. We presented a projection-free Lie algebra optimizer in PyTorch that is numerically stable and highly efficient. 

Furthermore, we provided a theoretical derivation showing that rotation and residual components must be scaled consistently with the interpolation parameter $\lambda$ to prevent representation collapse, which resolved a major coordinate-scale bottleneck in prior geometric merging works. Our empirical evaluation across a multi-task vision benchmark under clean and corrupted environments demonstrated that FWPA-FWF consistently, systematically, and substantially outperforms standard Unweighted OrthoMerge and Task Arithmetic, achieving absolute accuracy improvements of up to \textbf{+19.42\%} on clean data and \textbf{+31.78\%} under severe out-of-distribution blur.

\section*{Broader Impact Statement}
This work advances the field of model merging, enabling the creation of efficient, multi-task systems without the prohibitive carbon footprint and resource costs of joint pre-training. By providing a training-free framework that preserves expert capabilities under out-of-distribution corruptions, our method facilitates the deployment of robust machine learning models on edge devices and decentralized private networks. We do not foresee any direct negative societal impacts of this work, other than those general to the advancement of machine learning capabilities.

\newpage
\bibliography{submission}
\bibliographystyle{icml2026}

%%%%%%%%%%%%%%%%%%%%%%%%%%%%%%%%%%%%%%%%%%%%%%%%%%%%%%%%%%%%%%%%%%%%%%%%%%%%%%%
%%%%%%%%%%%%%%%%%%%%%%%%%%%%%%%%%%%%%%%%%%%%%%%%%%%%%%%%%%%%%%%%%%%%%%%%%%%%%%%
% APPENDIX
%%%%%%%%%%%%%%%%%%%%%%%%%%%%%%%%%%%%%%%%%%%%%%%%%%%%%%%%%%%%%%%%%%%%%%%%%%%%%%%
%%%%%%%%%%%%%%%%%%%%%%%%%%%%%%%%%%%%%%%%%%%%%%%%%%%%%%%%%%%%%%%%%%%%%%%%%%%%%%%
\newpage
\appendix
\onecolumn

\section{Detailed Experimental Setup and Hyperparameters}
\label{app:experiments}

To ensure maximum scientific reproducibility, we provide the exhaustive details of our experimental environment, training protocols, baseline parameterizations, and hyperparameters in this appendix.

\subsection{Dataset Details and Training Protocol}
Our multi-task vision benchmark consists of three 10-class datasets:
\begin{enumerate}
    \item \textbf{MNIST}: Handwritten digits dataset consisting of 60,000 training examples and 10,000 test examples, in grayscale and size $28 \times 28$.
    \item \textbf{FashionMNIST}: Grayscale clothing images consisting of 60,000 training examples and 10,000 test examples, size $28 \times 28$.
    \item \textbf{KMNIST (Kuzushiji-MNIST)}: Grayscale classical Japanese Hiragana character images consisting of 60,000 training examples and 10,000 test examples, size $28 \times 28$.
\end{enumerate}

For each dataset, the expert models are trained independently starting from a shared, randomly initialized base network encoder ($W_0$). We train the networks on their respective training sets for 3 epochs. We use the standard cross-entropy loss function. The optimizer is Adam with a learning rate of $1 \times 10^{-3}$ and a batch size of 128. We employ two distinct model architectures to evaluate the architectural generalizability of our methods:
\begin{enumerate}
    \item \textbf{CNN Encoder}: Contains two 2D convolutional layers. The first layer has 32 output channels with a kernel size of $3 \times 3$ and padding of 1. The second layer has 64 output channels with a kernel size of $3 \times 3$ and padding of 1. Both layers are followed by a ReLU activation function and a $2 \times 2$ MaxPool2D layer with stride 2. The output is flattened and passed through a fully connected projection layer to produce a 128-dimensional representation. Task-specific classification heads are single linear layers mapping from 128 dimensions to 10 output class logits.
    \item \textbf{MLP Encoder}: Contains three fully connected layers mapping from 784 dimensions (flattened $28 \times 28$ image) to 512, then to 256, and finally to a 128-dimensional output representation. Each layer is followed by a ReLU activation function. Task-specific classification heads are single linear layers mapping from 128 dimensions to 10 class logits.
\end{enumerate}

\subsection{Running Fisher Information Computation}
The diagonal running Fisher Information matrix $F_k$ is computed for each expert $k$ after independent fine-tuning. We estimate the Fisher Information diagonal using 1,000 samples randomly drawn from each expert's training set. For each sample, we execute a forward pass, compute the negative log-likelihood loss (cross-entropy), and accumulate the squared gradients with respect to the encoder parameters:
\begin{equation}
    F_k(p) = \frac{1}{N} \sum_{i=1}^N \left( \frac{\nabla \mathcal{L}(x_i, y_i)}{\partial W_k(p)} \right)^2
\end{equation}
To protect against numerical instabilities and prevent outlier parameters from completely dominating the objective, we apply a soft-bounding log-scale transformation:
\begin{equation}
    F'_k(p) = \log \left( 1 + \frac{F_k(p)}{\text{mean}(F_k) + \epsilon} \right)
\end{equation}
where $\epsilon = 10^{-8}$. As verified empirically in our unit testing suite, this transformation successfully bounds the maximum gradient weight to a stable range (e.g., contracting a range of $10^5$ down to a peak weight of less than $2.0$), stabilizing the Lie algebra optimization landscape.

\section{Detailed Baseline Implementations}
\label{app:baselines}

We systematically compared our method against four popular state-of-the-art model merging baselines. Each baseline is configured as follows:
\begin{enumerate}
    \item \textbf{Task Arithmetic (TA)} \cite{ilharco22}: We compute task vectors for each expert $k$ as $V_k = W_k - W_0$. The merged multi-task weights are computed linearly as $W_{TA} = W_0 + \lambda \sum_{k=1}^K V_k$ for various values of the interpolation scale $\lambda \in \{0.33, 0.40, 0.50\}$.
    \item \textbf{DARE} \cite{yu23}: We compute task vectors $V_k = W_k - W_0$. We randomly set a fraction $p_{drop} = 0.2$ of the task vector parameters to zero, and rescale the remaining parameters by $1 / (1 - p_{drop}) = 1.25$ to maintain structural scale. The scaled task vectors are then combined linearly via Task Arithmetic.
    \item \textbf{TIES-Merging} \cite{yadav23}: We compute task vectors $V_k = W_k - W_0$. We keep only the top 20\% magnitude parameters (keep\_rate = 0.2) in each task vector, setting the remaining 80\% to zero. We then construct a majority sign mask across the experts, and prune coordinates where the signs disagree with the majority. Finally, we average the non-conflicting task vectors and add them to $W_0$ scaled by $\lambda$.
    \item \textbf{Unweighted OrthoMerge} \cite{orthomerge}: We perform uniform Procrustes alignment on the convolutional filters/fully connected weights. For each expert, we compute the optimal rotation matrix $R_k \in \mathcal{O}(d_{out})$ that aligns $W_k$ with $W_0$ by solving $\min_R \|W_k - R W_0\|_F^2$ using Singular Value Decomposition (SVD), which gives $R_k = U V^T$ where $W_k W_0^T = U \Sigma V^T$. The Lie algebra matrices $A_k = \log(R_k)$ and residuals $E_k = W_k - R_k W_0$ are computed, merged by simple averaging, and reconstructed as $W_{merged} = \exp(A_{merged}) W_0 + E_{merged}$, where $A_{merged} = \frac{1}{K} \sum_k A_k$ and $E_{merged} = \lambda \sum_k E_k$.
\end{enumerate}

\section{Mathematical Background on Lie Algebra Optimization}
\label{app:lie_algebra}

In Section 2.3, we presented our projection-free Lie algebra optimizer. Here we provide the complete mathematical background justifying our skew-symmetric parameterization.

The special orthogonal group $SO(d)$ is the group of all $d \times d$ real matrices $R$ satisfying $R^T R = I$ and $\det(R) = +1$. The associated Lie algebra $\mathfrak{so}(d)$ consists of all $d \times d$ real skew-symmetric matrices $A$ satisfying $A^T = -A$. The exponential map $\exp: \mathfrak{so}(d) \to SO(d)$ is defined by the standard matrix exponential:
\begin{equation}
    \exp(A) = \sum_{n=0}^\infty \frac{A^n}{n!}
\end{equation}
Since any skew-symmetric matrix $A$ satisfies $A^T = -A$, we have:
\begin{equation}
    \left(\exp(A)\right)^T = \exp(A^T) = \exp(-A) = \left(\exp(A)\right)^{-1}
\end{equation}
which proves that $\exp(A)$ is strictly orthogonal. Additionally, because the Lie algebra is a connected vector space, the exponential of any skew-symmetric matrix yields a proper rotation (determinant $+1$).

To optimize the rotation matrix $R_k$ without requiring expensive projection steps (such as SVD or Gram-Schmidt orthonormalization) at every iteration, we represent $R_k$ as $\exp(B_k - B_k^T)$ where $B_k$ is an unconstrained real matrix of size $d_{out} \times d_{out}$. The gradient of our weighted Procrustes loss with respect to $B_k$ can be directly computed via standard backpropagation:
\begin{equation}
    \nabla_{B_k} \mathcal{L}_{opt} = \left(\frac{\partial \mathcal{L}_{opt}}{\partial R_k} \circ \frac{\partial \exp(A_k)}{\partial A_k}\right) \left( \frac{\partial(B_k - B_k^T)}{\partial B_k} \right)
\end{equation}
where $\circ$ represents tensor contractions of the exponential map's differential. Since PyTorch provides a fully differentiable, GPU-accelerated implementation of the matrix exponential via \texttt{torch.matrix\_exp}, we can optimize $B_k$ directly using standard gradient descent or Adam, completely avoiding unstable numerical matrix logarithm computations.

We use Adam with a learning rate of $\eta = 0.05$ for $T = 150$ optimization steps. The parameter matrix $B_k$ is initialized to $\mathbf{0}$, which corresponds to $A_k = \mathbf{0}$ and $R_k = I$.

\section{Definition of Out-of-Distribution Corruptions}
\label{app:corruptions}

To rigorously test model robustness and compatibility under out-of-distribution (OOD) test-time shifts, we evaluated our merged models across five distinct test-time environments:
\begin{enumerate}
    \item \textbf{Clean}: The original, uncorrupted test sets of MNIST, FashionMNIST, and KMNIST.
    \item \textbf{Gaussian Noise}: Adds independent, zero-mean Gaussian noise to each pixel. Let $x \in [0, 1]^{C \times H \times W}$ be the normalized input image. The corrupted image $x_{noise}$ is defined as:
    \begin{equation}
        x_{noise} = \text{clip}(x + \eta, 0, 1), \quad \text{where} \quad \eta \sim \mathcal{N}(0, \sigma^2 I)
    \end{equation}
    where we set the noise scale $\sigma = 0.2$.
    \item \textbf{Gaussian Blur}: Simulates camera out-of-focus blur by applying a 2D average pooling filter with a kernel size of $3 \times 3$, a stride of 1, and padding of 1:
    \begin{equation}
        x_{blur}(c, i, j) = \frac{1}{9} \sum_{m=-1}^1 \sum_{n=-1}^1 x(c, i+m, j+n)
    \end{equation}
    with zero-padding at the boundaries.
    \item \textbf{Contrast Reduction}: Simulates low-light contrast degradation by scaling all pixel intensities by a factor of 0.25:
    \begin{equation}
        x_{contrast} = 0.25 \times x
    \end{equation}
    reducing the dynamic range of the representations.
    \item \textbf{Image Rotation}: Simulates severe perspective and camera orientation changes by rotating the images clockwise by 90 degrees:
    \begin{equation}
        x_{rotation}(c, i, j) = x(c, H - 1 - j, i)
    \end{equation}
    This shift completely destroys standard translation symmetries and acts as a highly challenging OOD test.
\end{enumerate}

\section{Formal Theoretical Guarantees of Fisher-Weighted Procrustes Alignment}
\label{app:theory}

In this section, we provide a rigorous theoretical justification demonstrating why solving the Fisher-weighted orthogonal Procrustes problem mathematically minimizes the first-order and second-order approximations of task-specific performance degradation during model merging.

\subsection{Task Loss under Parameter Perturbation}
Let $\mathcal{L}_k: \mathbb{R}^D \to \mathbb{R}$ represent the training/evaluation loss of expert model $k$ for task $k$, where $D$ is the total number of flattened parameters in the shared encoder. Let $\theta_k \in \mathbb{R}^D$ represent the optimal parameter vector of expert $k$ obtained after independent task-specific fine-tuning.

Since $\theta_k$ is a local minimum of the loss function $\mathcal{L}_k$, the first-order necessary optimality condition implies that the gradient of the loss evaluated at $\theta_k$ is zero:
\begin{equation}
    \nabla \mathcal{L}_k(\theta_k) \approx \mathbf{0}
\end{equation}

When combining multiple expert models into a single multi-task system, the merged parameters $\theta_{merged}$ deviate from the individual optimal expert parameters $\theta_k$. Let $\Delta \theta_k = \theta_{merged} - \theta_k$ denote the parameter perturbation vector for expert $k$. Under a second-order Taylor series expansion of the loss function $\mathcal{L}_k$ around the expert weights $\theta_k$, we have:
\begin{equation}
    \mathcal{L}_k(\theta_{merged}) = \mathcal{L}_k(\theta_k + \Delta \theta_k) \approx \mathcal{L}_k(\theta_k) + \nabla \mathcal{L}_k(\theta_k)^T \Delta \theta_k + \frac{1}{2} \Delta \theta_k^T H_k(\theta_k) \Delta \theta_k
\end{equation}
where $H_k(\theta_k) \in \mathbb{R}^{D \times D}$ is the Hessian matrix of the loss $\mathcal{L}_k$ evaluated at $\theta_k$. 

Substituting the optimality condition $\nabla \mathcal{L}_k(\theta_k) \approx \mathbf{0}$, the task-specific performance degradation (or task interference) experienced by expert $k$ under the merged model is given by:
\begin{equation}
    \Delta \mathcal{L}_k = \mathcal{L}_k(\theta_{merged}) - \mathcal{L}_k(\theta_k) \approx \frac{1}{2} \Delta \theta_k^T H_k(\theta_k) \Delta \theta_k
\end{equation}

\subsection{Fisher Information Approximation of the Hessian}
In deep neural networks, computing and storing the full Hessian matrix $H_k(\theta_k)$ is computationally prohibitive due to its $D \times D$ dimensionality. Under standard probabilistic formulation where the loss function $\mathcal{L}_k$ is the negative log-likelihood of a conditional probability distribution, the expected Hessian at the optimal parameters equals the Fisher Information Matrix (FIM) $I(\theta_k)$:
\begin{equation}
    H_k(\theta_k) \approx I(\theta_k) = \mathbb{E}_{(x, y) \sim p_k} \left[ \nabla_\theta \log p_k(y | x; \theta_k) \nabla_\theta \log p_k(y | x; \theta_k)^T \right]
\end{equation}

To ensure computational tractability and scalability across millions of parameters, we employ a diagonal approximation of the Fisher Information Matrix, $F_k \approx \text{diag}(I(\theta_k))$. The performance degradation for expert $k$ can thus be written as a weighted sum of squared parameter differences:
\begin{equation}
    \Delta \mathcal{L}_k \approx \frac{1}{2} \sum_{p=1}^D F_k(p) \left( \theta_{merged}(p) - \theta_k(p) \right)^2 = \frac{1}{2} \| \sqrt{F_k} \odot (\theta_{merged} - \theta_k) \|_2^2
\end{equation}
where $\odot$ denotes the Hadamard (element-wise) product, and $F_k(p)$ is the running diagonal Fisher value of parameter $p$ for expert $k$.

\subsection{Minimizing Interference via Fisher-Weighted Alignment}
Our proposed Fisher-Weighted Procrustes Alignment (FWPA) algorithm decomposes the expert weights $W_k$ for each layer into an orthogonal rotation $R_k$ and a residual matrix $E_k$:
\begin{equation}
    W_k = R_k W_0 + E_k \implies E_k = W_k - R_k W_0
\end{equation}
For each expert, we optimize the rotation $R_k \in \mathcal{O}(d_{out})$ to minimize the Fisher-weighted Frobenius norm of the residual:
\begin{equation}
    \min_{R_k \in \mathcal{O}(d_{out})} \mathcal{J}(R_k) = \| \sqrt{F_k} \odot (W_k - R_k W_0) \|_F^2 = \| \sqrt{F_k} \odot E_k \|_F^2
\end{equation}

This optimization objective has a direct physical and geometric interpretation: it forces the residual error $E_k$ to concentrate in parameter directions where the Fisher information $F_k$ is extremely small (i.e., highly flat/redundant regions of the loss landscape). Conversely, in directions where the parameters are highly sensitive (high-Fisher), the residual $E_k$ is forced to be extremely close to zero, meaning that the coordinate alignment under $R_k$ preserves the original expert weights perfectly.

When merging the experts, our Fisher-Weighted Fusion (FWF) scheme combines the residuals and Lie algebra rotations using parameter-wise Fisher weights:
\begin{equation}
    E_{merged} = \lambda \sum_{k=1}^K \left( K w_k(p) \right) E_k, \quad \text{where} \quad w_k(p) = \frac{F_k(p)}{\sum_{j=1}^K F_j(p) + \epsilon}
\end{equation}

By choosing parameter-wise Fisher weights $w_k(p)$ that are proportional to the relative task sensitivity of each expert, the merged residual $E_{merged}(p)$ at each parameter $p$ is dominated by the expert that cares most about that parameter. Specifically, if expert $k$ has a highly sensitive parameter at coordinate $p$ (large $F_k(p)$) while other experts have low sensitivity, the weight $w_k(p) \approx 1$ and $w_j(p) \approx 0$ for $j \neq k$. This decouples the Euclidean parameter space and prevents task interference.

Combining the optimal rotation alignment with Fisher-weighted fusion guarantees that:
\begin{equation}
    \sum_{k=1}^K \Delta \mathcal{L}_k \approx \frac{1}{2} \sum_{k=1}^K \| \sqrt{F_k} \odot (W_{merged} - W_k) \|_F^2
\end{equation}
is mathematically minimized across the multi-task parameter space. This provides a rigorous theoretical foundation explaining why FWPA-FWF achieves outstanding performance gains and out-of-distribution robustness compared to unweighted coordinate alignment and standard linear merging.

\section{Full Hyperparameter Reference Table}
\label{app:table}

For convenience and direct reference, we summarize all the hyperparameter values used in our experts training, Fisher computation, baseline setups, and Lie algebra optimization phases in Table \ref{tab:hyperparameters}.

\begin{table}[h]
\caption{Complete hyperparameter reference table for all phases of our model merging and evaluation workflow.}
\label{tab:hyperparameters}
\vskip 0.15in
\begin{center}
\begin{small}
\begin{sc}
\begin{tabular}{llr}
\toprule
Phase & Hyperparameter & Value \\
\midrule
Expert Training & Optimizer & Adam \\
& Learning Rate & $1 \times 10^{-3}$ \\
& Batch Size & 128 \\
& Epochs & 3 \\
& Loss Function & Cross-Entropy \\
\midrule
Fisher Estimation & Number of Samples & 1,000 \\
& Batch Size & 1 \\
& Numerical Stabilizer $\epsilon$ & $10^{-8}$ \\
\midrule
DARE Baseline & Parameter Drop Rate ($p_{drop}$) & 0.2 \\
TIES-Merging Baseline & Parameter Keep Rate ($keep\_rate$) & 0.2 \\
\midrule
Lie Algebra Optimization & Optimizer & Adam \\
& Learning Rate ($\eta$) & 0.05 \\
& Total Optimization Steps ($T$) & 150 \\
& Initialization matrix ($B_k$) & $\mathbf{0}$ \\
\midrule
Model Merging & Interpolation Scale ($\lambda$) & 0.33, 0.40, 0.50 \\
& Number of Experts ($K$) & 3 \\
\bottomrule
\end{tabular}
\end{sc}
\end{small}
\end{center}
\vskip -0.1in
\end{table}

\end{document}
